%!TeX program = lualatex
\documentclass[10pt, aspectratio=169, t]{beamer}
\include{preamble.tex}

% TITLE

\title{\textbf{Hydraulic Enceladus}}
\subtitle{COSC-P project}
\author{Kamil Belán}
\date{\today}
\addtobeamertemplate{frametitle}{}{\vspace{0.5em}}
\begin{document}

% TITLE SLIDE 
{
	\setbeamercolor{background canvas}{bg=PaperBg}
	\begin{frame}[plain]
		\vspace{2cm}
		{\Huge \color{DragonRed} \inserttitle} \\
		\vspace{0.3cm}
		{\Large \color{DragonGray} \insertsubtitle} \\
		\vspace{1.5cm}
		{\insertauthor}
	\end{frame}
}

\section{Motivation}
\begin{frame}{Motivation}

\end{frame}

\section{Problem settings}
\begin{frame}{Modelling choices}
	\begin{itemize}
		\item the crevasses are deep narrow \emph{effectively} rectangular domains,
			\vfill
		\item the fluid is water, \textit{i.e.} \textbf{incompressible} Newtonian\footnote{Much higher viscosity has been used in the simulation.}, $\rho = \SI{1000}{\kilogram /\meter^3}, \mu =\SI{1,5}{\milli\pascal  \second}$,
			\vfill
		\item the only body force is \textbf{gravity} $\rho \vb{g}, g = \SI{0.113}{\meter / \second^2}, $
			\vfill
		\item the fluid can \emph{freely slip} along the walls,
			\vfill
		\item the motion of the walls is \emph{harmonic}
			\vfill
	\end{itemize}
\end{frame}

\begin{frame}{Governing equations}
	We are solving for $\left(\vb{v}, p, \vb{u}\right)$ with $\Omega_0 = (0, 1) \times (0, 10)$ given such that $\forall t \in (0, T), T > 0$ inside the domain
	\begin{math-block}{Governing equations}
		\begin{equations}[single,*]
			\Omega_t = \Big\{\vb{x} + \vb{u}(t, \vb{x}) \mid \vb{x} \in \Omega_0\Big\} \subset \mathbb{R}^2
		\end{equations}
		it holds $\forall t \in (0, T)$
		\begin{equations}[columns,*]
			\bigtriangleup \vb{u} &= 0, \\
			\rho\left(\partial_t \vb{v} + \left(\vb{v} - \partial_t \vb{u}\right) \vdot \grad \vb{v}\right) &= \divergence{\symbb{T}} + \rho \vb{g}, \\ 
			\divergence{\vb{v}} &= 0, \\
			\symbb{T} &= -p \symbb{I} + 2 \mu \symbb{D}.
		\end{equations}
	\end{math-block}
\end{frame}

\begin{frame}{Boundary conditions}
	Denote $\partial \Omega_t = \Gamma_{\, \text{bot} \,} \cup \Gamma_{\, \text{left} \,} \cup \Gamma_{\, \text{right} \,} \cup \Gamma_{\, \text{top} \,}.$ The trio $\left(\vb{v}, p, \vb{u} \right)$ must satisfy on $\left(0, T \right) \times \partial \Omega_t$

	\begin{math-block}{Boundary conditions}
		\begin{equations}[columns,*]
			u_x &= u_{0} \sin(2 \pi f t), \partial_{x} u_y = 0, \, \text{on} \, \Gamma_{\, \text{left} \,}, \\
			u_x &= -u_{0} \sin(2 \pi f t), \partial_{x} u_y = 0\, \text{on} \, \Gamma_{\, \text{right} \,}, \\
			u_y &= 0, \partial_{y} u_x = 0\, \text{on} \, \Gamma_{\, \text{bot} \,}, \\
			v_x &= 2 u_0 \pi f \cos(2 \pi f t), \partial_{x} v_y = 0, \, \text{on} \, \Gamma_{\, \text{left} \,}, \\
			v_x &= -2 u_0 \pi f \cos(2 \pi f t), \partial_{x} v_y = 0, \, \text{on} \, \Gamma_{\, \text{right} \,} ,\\
			v_y &= 0, \partial_{y} v_x = 0 \, \text{on} \, \Gamma_{\, \text{bot} \,}, \\
			\partial_t \vb{u} &= \vb{v}, \, \text{on} \, \Gamma_{\, \text{top} \,}, \\
			\symbb{T} \vb{n} &= \vb{0}, \, \text{on} \, \Gamma_{\, \text{top} \,}.
		\end{equations}
	\end{math-block}
\end{frame}

\section{Numerical solution}
\begin{frame}{ Implementation strategy}
	To solve the problem numerically, we use the \emph{FEniCS} library with the choices:

	\vfill	

	\begin{itemize}
		\item \textbf{Discretization:} $\left(P_2, P_1, P_2\right)$ elements for $\left(\vb{v}, p, \vb{u}\right)$ \textit{i.e.} Taylor-Hood for fluid, $P_2$ for mesh,
			\vfill
		\item \textbf{Coupling}: full ALE method with the monolithic formulation for $(\vb{v}, p, \vb{u})$.
			\vfill
		\item \textbf{Kinematic top boundary condition}: (symmetric) Nitsche Method for weak enforcement of the free surface BC
			\vfill
		\item \textbf{Time Integration:} implicit Crank-Nicolson $\theta$-scheme
			\vfill
		\item \textbf{Solver:} Newton-based with MUMPS direct solver.
			\vfill
	\end{itemize}
\end{frame}

\begin{frame}{Parameters}
	The following parameters have been chosen to demonstrate the \textbf{qualitative behaviour} of the model

	\vfill

	\begin{itemize}
		\item \emph{time step}: $dt = 0.05$ 
			\vfill
		\item \emph{mesh resolution}: 10 cells in horizontal direction and 400 cells in vertical direction
			\vfill
		\item \emph{frequency of wall motion}: $f =\SI{0.05}{\second^{-1}}, $
			\vfill
		\item \emph{amplitude of wall motion}: $u_0 = \SI{0.40}{\meter}, $
			\vfill
		\item \emph{gravitational acceleration}: $g = \SI{0.113}{\meter /\second^2}, $
			\vfill
		\item \emph{density}: $\rho = \SI{1000}{\kilogram / \meter^3}, $
			\vfill
		\item \emph{dynamic viscosity}: $\mu = \SI{1490}{\Pa \second}, $
	\end{itemize}
\end{frame}

\section{Results}

\begin{frame}{Close view}

\end{frame}

\begin{frame}{Complete view}
    \centering
    % Adjust 'width' to fill the frame nicely
    \movie[height=0.9\textheight, width=0.9\textwidth, poster, showcontrols]{}{zoom_out.mp4}
    \href{run:./videos/zoom_out.mp4?autostart&loop}{
        \includegraphics[height=0.85\textheight, keepaspectratio]{figures/zoom_out.png}
    }
    \vspace{0.5em}
    \footnotesize{\textit{Cyclic squeezing motion of the Enceladus crevasse.}}
\end{frame}


\end{document}
